% Options for packages loaded elsewhere
\PassOptionsToPackage{unicode}{hyperref}
\PassOptionsToPackage{hyphens}{url}
\documentclass[
  ignorenonframetext,
]{beamer}
\newif\ifbibliography
\usepackage{pgfpages}
\setbeamertemplate{caption}[numbered]
\setbeamertemplate{caption label separator}{: }
\setbeamercolor{caption name}{fg=normal text.fg}
\beamertemplatenavigationsymbolsempty
% remove section numbering
\setbeamertemplate{part page}{
  \centering
  \begin{beamercolorbox}[sep=16pt,center]{part title}
    \usebeamerfont{part title}\insertpart\par
  \end{beamercolorbox}
}
\setbeamertemplate{section page}{
  \centering
  \begin{beamercolorbox}[sep=12pt,center]{section title}
    \usebeamerfont{section title}\insertsection\par
  \end{beamercolorbox}
}
\setbeamertemplate{subsection page}{
  \centering
  \begin{beamercolorbox}[sep=8pt,center]{subsection title}
    \usebeamerfont{subsection title}\insertsubsection\par
  \end{beamercolorbox}
}
% Prevent slide breaks in the middle of a paragraph
\widowpenalties 1 10000
\raggedbottom
\AtBeginPart{
  \frame{\partpage}
}
\AtBeginSection{
  \ifbibliography
  \else
    \frame{\sectionpage}
  \fi
}
\AtBeginSubsection{
  \frame{\subsectionpage}
}
\usepackage{iftex}
\ifPDFTeX
  \usepackage[T1]{fontenc}
  \usepackage[utf8]{inputenc}
  \usepackage{textcomp} % provide euro and other symbols
\else % if luatex or xetex
  \usepackage{unicode-math} % this also loads fontspec
  \defaultfontfeatures{Scale=MatchLowercase}
  \defaultfontfeatures[\rmfamily]{Ligatures=TeX,Scale=1}
\fi
\usepackage{lmodern}
\ifPDFTeX\else
  % xetex/luatex font selection
\fi
% Use upquote if available, for straight quotes in verbatim environments
\IfFileExists{upquote.sty}{\usepackage{upquote}}{}
\IfFileExists{microtype.sty}{% use microtype if available
  \usepackage[]{microtype}
  \UseMicrotypeSet[protrusion]{basicmath} % disable protrusion for tt fonts
}{}
\makeatletter
\@ifundefined{KOMAClassName}{% if non-KOMA class
  \IfFileExists{parskip.sty}{%
    \usepackage{parskip}
  }{% else
    \setlength{\parindent}{0pt}
    \setlength{\parskip}{6pt plus 2pt minus 1pt}}
}{% if KOMA class
  \KOMAoptions{parskip=half}}
\makeatother
\usepackage{longtable,booktabs,array}
\newcounter{none} % for unnumbered tables
\usepackage{calc} % for calculating minipage widths
\usepackage{caption}
% Make caption package work with longtable
\makeatletter
\def\fnum@table{\tablename~\thetable}
\makeatother
\setlength{\emergencystretch}{3em} % prevent overfull lines
\providecommand{\tightlist}{%
  \setlength{\itemsep}{0pt}\setlength{\parskip}{0pt}}
% Slide layout defaults for warning-clean scientific decks.
\usepackage{etoolbox}
\IfFileExists{xurl.sty}{\usepackage{xurl}}{}
\IfFileExists{seqsplit.sty}{\usepackage{seqsplit}}{\newcommand{\seqsplit}[1]{#1}}
\protected\def\breakseq#1{\seqsplit{#1}}
\protected\def\breaktt#1{\begingroup\ttfamily\seqsplit{#1}\endgroup}
\setlength{\emergencystretch}{6em}
\tolerance=5000
\hbadness=10000
\hfuzz=1pt
\setlength{\tabcolsep}{2pt}
\AtBeginEnvironment{longtable}{\tiny\renewcommand{\arraystretch}{0.86}\setlength{\tabcolsep}{1pt}}
\AtBeginEnvironment{tabular}{\tiny\renewcommand{\arraystretch}{0.86}\setlength{\tabcolsep}{1pt}}

% Natbib and cross-reference fallbacks — slides don't load natbib
% or cleveref, but combined-PDF manuscript prose may emit these
% commands. The fallback renders citations as a bracketed key list
% and cross-references as detokenized labels so slides don't fail on
% undefined control sequences or raw underscores. \providecommand is
% a no-op if packages load later.
\providecommand{\citep}[1]{[#1]}
\providecommand{\citet}[1]{#1}
\providecommand{\citealp}[1]{#1}
\providecommand{\citeauthor}[1]{#1}
\providecommand{\citeyear}[1]{#1}
\providecommand{\cref}[1]{\texttt{\detokenize{#1}}}
\providecommand{\Cref}[1]{\texttt{\detokenize{#1}}}
\usepackage{bookmark}
\IfFileExists{xurl.sty}{\usepackage{xurl}}{} % add URL line breaks if available
\urlstyle{same}
% fallback for those not using the hyperref driver hyperxmp:
\makeatletter
\@ifundefined{xmpquote}{\newcommand{\xmpquote}[1]{#1}}{}
\makeatother
\hypersetup{
  hidelinks,
  pdfcreator={LaTeX via pandoc}}

\author{\texorpdfstring{}{}}
\date{}

\begin{document}

\section{Introduction}\label{introduction}

\begin{frame}[allowframebreaks]{Research Tools as Epistemic
Infrastructure}
\protect\phantomsection\label{research-tools-as-epistemic-infrastructure}
Scientific research operates through a layered ecology of tools,
documents, and practices---each shaping what can be known, communicated,
and verified. When these layers are fragmented, the artifacts of
research (manuscripts, data, code, figures) drift out of alignment with
one another, creating gaps that are structural rather than incidental.
The ``reproducibility crisis'' is one symptom of this deeper
misalignment: a 2016 \emph{Nature} survey of 1,576 researchers found
that 70\% had tried and failed to reproduce another scientist's
experiments, and more than half had failed to reproduce their own
\breakseq{[@baker2016reproducibility]}. Freedman et al.~estimate that the
biomedical industry alone loses \$28 billion annually to irreproducible
preclinical research \breakseq{[@freedman2015economics]}. But reproducibility
is only the most visible face of a broader problem. Research software
engineering, epistemic integrity, and the coordination between human
researchers and AI collaborators all depend on the same underlying
question: whether the tools that produce research artifacts can
themselves be made transparent, testable, and self-documenting. Nosek et
al. \breakseq{[@nosek2018preregistration]} have argued that the preregistration
revolution---requiring researchers to commit to analytical plans before
data collection---is a necessary structural reform; we extend this logic
to the entire research build pipeline.

One root cause is fragmentation---of attention (in the mind) as well as
in the cyberphysical niche (documents, versions, and reproducibility
artefacts). A typical research project scatters its artifacts across
disconnected tools: LaTeX editors \breakseq{[@lamport1994latex]} for prose,
Jupyter notebooks for analysis, ad-hoc shell scripts for figure
generation, and manual copy-paste for integrating results into
manuscripts. Each boundary between tools is a potential locus of
desynchronization. The version of the figure embedded in the PDF may not
match the version of the code that ostensibly generated it. The test
suite, if it exists at all, likely tests the code in isolation from the
rendering pipeline. Pimentel et al. \breakseq{[@pimentel2019jupyter]} analyzed
1.4 million Jupyter notebooks from GitHub and found that only 24\% could
be successfully re-executed, with 36\% producing different
results---quantifying the reproducibility cost of notebook-based
workflows. Peng \breakseq{[@peng2011reproducible]} argues that reproducibility
in computational science requires, at minimum, that the data and code
underlying a published result be available for independent
verification---yet the tools for enforcing this standard remain ad hoc.
Indeed, even the terminology is fractured: Barba
\breakseq{[@barba2018terminologies]} documents how ``reproducibility,''
``replicability,'' and ``repeatability'' carry conflicting definitions
across disciplines, undermining cross-field standards.
\end{frame}

\begin{frame}[fragile,allowframebreaks]{Related Work}
\protect\phantomsection\label{related-work}
Gentleman and Temple Lang \breakseq{[@gentleman2007research]} introduced the
concept of a \emph{research compendium}---a single unit of scholarly
communication bundling code, data, and narrative. This vision has driven
two decades of tooling, which can be grouped into four categories:
workflow managers, literate programming systems, containerization
approaches, and best-practice frameworks.

\begin{block}{Workflow Managers}
\protect\phantomsection\label{workflow-managers}
\textbf{Snakemake} \breakseq{[@koster2012snakemake]} uses a rule-based,
Python-derived DSL to specify computational workflows as directed
acyclic graphs of file-producing steps. It supports containerized
execution via Conda and Singularity environments. Snakemake 9.x
(2024--2025) introduced a plugin architecture for extended execution
backends and storage providers, yet its scope remains computational
pipeline orchestration---it does not integrate manuscript rendering,
testing enforcement, or provenance watermarking.

\textbf{Nextflow} \breakseq{[@ditommaso2017nextflow]} employs a dataflow
programming paradigm with native support for container-based execution
across heterogeneous computing environments (local, SLURM, AWS). Like
Snakemake, Nextflow excels at bioinformatics pipeline parallelism but
does not address manuscript production, document integrity, or the
testing--publication coupling that characterizes research
reproducibility.

\textbf{CWL} (Common Workflow Language) \breakseq{[@amstutz2016cwl]} provides a
portable, YAML-based standard for describing computational workflows and
their dependencies. Its strength lies in interoperability across
execution engines (cwltool, Toil, Arvados), but it requires external
tooling for manuscript generation and offers no built-in testing or
provenance framework.
\end{block}

\begin{block}{Literate Programming and Publication Tools}
\protect\phantomsection\label{literate-programming-and-publication-tools}
Knuth's literate programming \breakseq{[@knuth1984literate]} established the
principle that programs should be authored as documents intended for
human comprehension. Schulte et al. \breakseq{[@schulte2012multilanguage]}
extended this to multi-language computing environments (Org-mode),
demonstrating that literate programming could span languages and output
formats.

\textbf{Quarto} \breakseq{[@allaire2024quarto]} extends the R Markdown
tradition to support Python, Julia, and Observable, rendering to PDF,
HTML, and Word. Quarto integrates code execution with document
rendering, achieving a modern form of literate programming, but it does
not enforce testing thresholds, manage multi-project repositories, or
provide cryptographic provenance.

\textbf{Jupyter Book} \breakseq{[@kluyver2016jupyter]} builds on Jupyter
notebooks to produce publication-quality documents via Sphinx. While
powerful for interactive exploration, Jupyter's notebook format
introduces execution-order fragility \breakseq{[@pimentel2019jupyter]} and does
not naturally support the separation of logic from orchestration that
characterizes maintainable research software.

\textbf{R Markdown} \breakseq{[@xie2018dynamic]} pioneered knitr-based dynamic
documents that weave code and prose. Its ecosystem is rich but
R-centric, and it lacks the multi-project management, infrastructure
testing, and provenance embedding that characterize \texttt{template/}.

\textbf{Typst} \breakseq{[@madje2023typst]} is an emerging markup-based
typesetting system with incremental compilation and a programmable
scripting layer. While Typst offers faster compilation than LaTeX and a
more modern authoring experience, it does not integrate testing,
provenance, or multi-project pipeline management.
\end{block}

\begin{block}{Containerization and Environment Capture}
\protect\phantomsection\label{containerization-and-environment-capture}
\textbf{Docker} \breakseq{[@boettiger2015docker]} addresses reproducibility at
the environment level---packaging operating system, libraries, and code
into portable containers. While Docker solves the ``works on my
machine'' problem, containerization is complementary to, not a
replacement for, the architectural concerns addressed here: Docker does
not enforce testing, embed provenance, or manage multi-project
manuscript workflows.
\end{block}

\begin{block}{Best-Practice Frameworks and Data Standards}
\protect\phantomsection\label{best-practice-frameworks-and-data-standards}
Wilson et al. \breakseq{[@wilson2017good]} define ``good enough'' practices for
scientific computing, emphasizing version control, testing, and
documentation. Sandve et al. \breakseq{[@sandve2013ten]} propose ten rules for
reproducible computational research. Piccolo and Frampton
\breakseq{[@piccolo2016tools]} systematically survey tools for computational
reproducibility, finding that environment isolation, workflow
automation, and documentation generation address complementary but
non-overlapping reproducibility concerns---yet no single tool unifies
all three. Stodden et al. \breakseq{[@stodden2016enhancing]} advocate for
enhanced computational method transparency. The FAIR principles
\breakseq{[@wilkinson2016fair]}---Findable, Accessible, Interoperable,
Reusable---establish a standard for data stewardship that has been
widely adopted by funding agencies and journals. Lamprecht et al.
\breakseq{[@lamprecht2020towards]} formalize ``Towards FAIR Principles for
Research Software,'' providing the conceptual scaffolding that Barker et
al. \breakseq{[@barker2022fair4rs]} would soon operationalize as the FAIR4RS
initiative---recognizing that software has execution, composability, and
dependency-management requirements that data-centric FAIR does not
address. Cohen et al. \breakseq{[@cohen2021four]} characterize the four pillars
of research software engineering (software sustainability, software
quality, community building, and policy advocacy), situating formal
testing and provenance practices within a broader RSE governance
framework. Garijo et al. \breakseq{[@garijo2024fairsoft]} operationalize
FAIR4RS through the FAIRsoft evaluator, an automated assessment
framework that scores research software against 17+ quality indicators
including executability, metadata richness, and documentation
completeness. Goble et al. \breakseq{[@goble2020fair]} extend FAIR to
computational workflows specifically, arguing that workflow provenance
requires first-class treatment in scientific computing infrastructure.
Nüst et al. \breakseq{[@nust2017containerization]} introduce the
\emph{executable research compendium} (ERC), extending Gentleman and
Temple Lang's compendium concept with containerized, interactive
reproduction environments. The W3C PROV data model
\breakseq{[@moreau2013provdm]} provides a formal vocabulary for expressing
provenance records, while in-toto \breakseq{[@torresarias2019intoto]} provides
a framework for end-to-end software supply chain integrity verification,
and SLSA \breakseq{[@openssf2023slsa]} (Supply-chain Levels for Software
Artifacts) extends this to graduated, attestation-based supply-chain
security levels for build pipelines. These frameworks articulate
\emph{what} reproducible research requires but do not provide an
integrated \emph{how}---they lack the tooling, enforcement mechanisms,
and architectural patterns that translate standards into practice.
\end{block}

\begin{block}{The Gap}
\protect\phantomsection\label{the-gap}
Despite advances in FAIR4RS principles \breakseq{[@barker2022fair4rs;
@lamprecht2020towards]}, automated FAIR software assessment
\breakseq{[@garijo2024fairsoft]}, FAIR computational workflow standards
\breakseq{[@goble2020fair]}, supply-chain attestation frameworks
\breakseq{[@torresarias2019intoto; @openssf2023slsa]}, and the preregistration
revolution \breakseq{[@nosek2018preregistration]}, no existing system
integrates six cross-cutting concerns into a single enforced pipeline:
(1) end-to-end pipeline orchestration with testing enforcement, (2)
multi-format manuscript rendering, (3) cryptographic provenance
embedding, (4) multi-project repository management, (5) FAIR-aligned
software stewardship, and (6) AI-agent collaboration via structured
documentation. Each existing framework addresses a subset; none provides
the unified enforcement mechanism. The detailed tool-by-tool comparison
is developed in the
\href{05c_comparison.md\#comparison-to-existing-tools}{Comparison to
Existing Tools} section; the summary table below captures the gap
landscape.
\end{block}

\begin{block}{Summary: Gaps Left by Existing Tools}
\protect\phantomsection\label{summary-gaps-left-by-existing-tools}
The six concerns identified above map onto existing tool categories as
follows:

{\def\LTcaptype{none} % do not increment counter
\begin{longtable}[]{@{}
  >{\raggedright\arraybackslash}p{(\linewidth - 4\tabcolsep) * \real{0.2222}}
  >{\raggedright\arraybackslash}p{(\linewidth - 4\tabcolsep) * \real{0.4000}}
  >{\raggedright\arraybackslash}p{(\linewidth - 4\tabcolsep) * \real{0.3778}}@{}}
\toprule\noalign{}
\begin{minipage}[b]{\linewidth}\raggedright
Gap
\end{minipage} & \begin{minipage}[b]{\linewidth}\raggedright
Partially Addressed By
\end{minipage} & \begin{minipage}[b]{\linewidth}\raggedright
Not Addressed By
\end{minipage} \\
\midrule\noalign{}
\endhead
Pipeline orchestration & Snakemake 9.x, Nextflow 25.x, CWL 1.2 & Quarto,
Jupyter Book, R Markdown, Typst, Overleaf, Prism \\
Manuscript rendering & Quarto 1.x, Jupyter Book 2.x, R Markdown, Typst,
Overleaf (2025), Prism & Snakemake, Nextflow, CWL, DVC \\
Testing enforcement & --- & All existing tools \\
Cryptographic provenance & SLSA (build-level attestation only) & All
research-focused tools \\
Multi-project management & --- & All existing tools \\
AI-agent documentation & Overleaf (partial co-author AI), Prism (partial
context reasoning) & All pipeline/workflow tools \\
Agentic skill protocol (MCP-aligned) & --- & All existing tools \\
\bottomrule\noalign{}
\end{longtable}
}

No existing system addresses all six concerns within a single enforced
pipeline.
\end{block}
\end{frame}

\begin{frame}[fragile,allowframebreaks]{\texttt{template/}: An
Integrated Solution}
\protect\phantomsection\label{template-an-integrated-solution}
\texttt{template/} was conceived as a structural antidote to this
fragmentation. Rather than adding reproducibility as an afterthought---a
Docker container wrapping an already-disjointed workflow
\breakseq{[@boettiger2015docker]}---the template enforces integrity at the
architectural level. It realizes Gentleman and Temple Lang's research
compendium vision \breakseq{[@gentleman2007research]} at repository scale,
bundling code, data, tests, manuscripts, and provenance into a single,
pipeline-enforced system with version-controlled infrastructure
\breakseq{[@ram2013git]}. It stands on four primary pillars:

\begin{enumerate}
\item
  \textbf{Ergonomic Modularity}: A Two-Layer Architecture cleanly
  separates globally shared infrastructure (logging, rendering,
  validation, steganography) from project-specific logic (manuscripts,
  scripts, data). 23 infrastructure subdirectories (20 importable
  packages) comprising \textasciitilde567 Python modules provide
  reusable services; projects consume them without modification.
\item
  \textbf{Execution Integrity}: A Zero-Mock testing policy where
  pipeline advancement is contingent on test passage. Infrastructure
  tests must achieve 60\% coverage; project tests must achieve 90\%. All
  tests use real filesystem operations, real subprocess calls, and real
  network connections---no mock objects, no fake services, no synthetic
  test doubles. \textasciitilde7,375 infrastructure tests and 1236+
  project tests enforce this standard.
\item
  \textbf{Automated Provenance}: Steganographic watermarking and
  cryptographic hashing are integrated directly into the rendering
  pipeline. Every generated PDF carries a SHA-256 fingerprint, an
  alpha-channel text overlay encoding the build timestamp and commit
  hash, and optionally a QR code linking to the repository. Provenance
  is not asserted by policy; it is enforced by architecture.
\item
  \textbf{AI-Agent Collaboration and Skill-Based Agentic Operations}: A
  three-tier documentation architecture enables AI agents to operate at
  every level of the system. At the system level, \texttt{CLAUDE.md}
  asserts global architectural constraints. At the structural level,
  \texttt{AGENTS.md} files at every directory expose local API surfaces,
  file inventories, and integration contracts. At the module level,
  \texttt{SKILL.md} files---written to a discoverable YAML+Markdown
  schema aligned with the Model Context Protocol
  \breakseq{[@anthropic2024mcp]}---define each infrastructure module as a
  reusable, self-describing tool. An agent invoking
  \breaktt{infrastructure.rendering} does not need to read source code:
  it reads the \breaktt{rendering/SKILL.md}, which declares the module's
  name, description, key imports, and example invocations in a
  machine-parseable YAML frontmatter block. This architecture is the
  practical realization of the skill-library paradigm established in the
  agent literature: Yao et al.'s ReAct framework \breakseq{[@yao2023react]}
  demonstrated that interleaving reasoning traces with tool invocations
  dramatically improves LLM reliability; Schick et al.'s Toolformer
  \breakseq{[@schick2023toolformer]} showed that self-supervised tool use can
  be bootstrapped from natural language; Wang et al.'s Voyager
  \breakseq{[@wang2023voyager]} proved that growing skill libraries enable
  open-ended autonomous exploration in complex environments.
  \texttt{template/} instantiates this vision in the domain of
  scientific research infrastructure: each \texttt{SKILL.md} is a
  Voyager-style skill, each pipeline stage is a ReAct action, and the
  full infrastructure layer constitutes a composable, protocol-aligned
  skill library for scientific computation.
\end{enumerate}
\end{frame}

\begin{frame}[fragile,allowframebreaks]{Scope and Contributions}
\protect\phantomsection\label{scope-and-contributions}
This paper is itself a product of the template it describes. The metrics
populating its tables were computed by the introspection module
documented in the \href{03a_architecture.md\#methods}{Methods}; the
figures were rendered by the visualization code validated by the test
suite described in \href{03e_quality.md\#quality-assurance}{Quality
Assurance}; the PDF carrying these words was assembled by the same
YAML-declared pipeline whose architecture is the subject of the
\href{04_results.md\#results}{Results}. This self-productive
loop---where the system that is described is also the system that
produces the description---is not incidental but structural, a concrete
demonstration that \texttt{template/} can sustain the full lifecycle
from source code to published artifact within a single,
version-controlled, pipeline-enforced repository. Our contributions are:

\begin{itemize}
\tightlist
\item
  A formal description of the Two-Layer Architecture and Standalone
  Project Paradigm that enables N independent research projects to share
  infrastructure without coupling.
\item
  A detailed specification of the 12-stage DAG in
  \breaktt{infrastructure/core/pipeline/pipeline.yaml}: default full runs
  use 10 stages; \texttt{-\/-core-only} runs 8; opt-in bundle and
  archival stages via \texttt{-\/-tags}.
\item
  A comparative analysis positioning \texttt{template/} against nine
  peer tools---Snakemake, Nextflow, CWL, Quarto, Jupyter Book, R
  Markdown, DVC, Overleaf, OpenAI Prism---across fourteen feature
  dimensions, demonstrating that \texttt{template/} uniquely bundles the
  fourteen enforcement capabilities enumerated in §Results.
\item
  An empirical evaluation across the canonical exemplar projects under
  \breaktt{projects/templates/}
  (\breaktt{templates/template\_active\_inference},
  \breaktt{templates/template\_autoresearch\_project},
  \breaktt{templates/template\_autoscientists},
  \breaktt{templates/template\_code\_project},
  \breaktt{templates/template\_madlib},
  \breaktt{templates/template\_newspaper},
  \breaktt{templates/template\_prose\_project},
  \breaktt{templates/template\_sia},
  \breaktt{templates/template\_template},
  \breaktt{templates/template\_textbook})---including this meta
  manuscript from \breaktt{projects/templates/template\_template}, which
  exercises introspection-derived metrics.
\item
  A security analysis of the steganographic provenance layer, including
  a formal threat model and tamper-detection capabilities aligned with
  the W3C PROV data model \breakseq{[@moreau2013provdm]} and SLSA
  \breakseq{[@openssf2023slsa]}.
\item
  An open-source reference implementation available at
  \breaktt{github.com/docxology/template}.
\end{itemize}
\end{frame}

\begin{frame}[allowframebreaks]{Paper Organization}
\protect\phantomsection\label{paper-organization}
The \href{03a_architecture.md\#methods}{Methods} describe the Two-Layer
Architecture, Thin Orchestrator pattern, pipeline stages, and AI
collaboration model. \href{04_results.md\#results}{Results} present
quantitative metrics from multi-project execution, coverage analysis,
and steganographic benchmarks. The
\href{05a_zeromock_tradeoff.md\#the-zero-mock-tradeoff}{Discussion}
addresses the Zero-Mock tradeoff, scalability implications, a detailed
tool comparison, and future directions. The
\href{06_infrastructure_modules.md\#infrastructure-module-reference}{Infrastructure
Module Reference} provides detailed documentation for all 23
subdirectories.
\href{07_security_provenance.md\#security-and-provenance}{Security and
Provenance} describes the steganographic and cryptographic integrity
layer. The
\href{08a_appendix_pipeline.md\#appendix-pipeline}{Appendices} provide
pipeline, configuration, and comparative references.
\end{frame}

\end{document}
